        %%******************************************%%
        %%                                          %%
        %%        Modello di tesi di laurea         %%
        %%            di Andrea Giraldin            %%
        %%                                          %%
        %%             2 novembre 2012              %%
        %%                                          %%
        %%******************************************%%

% I seguenti commenti speciali impostano:
% 1. 
% 2. PDFLaTeX come motore di composizione;
% 3. tesi.tex come documento principale;
% 4. il controllo ortografico italiano per l'editor.

% !TEX encoding = UTF-8
% !TEX TS-program = pdflatex
% !TEX root = tesi.tex
% !TEX spellcheck = it-IT

\documentclass[10pt,                    					% corpo del font principale
               a4paper,                 					% carta A4
               twoside,                 					% impagina per fronte-retro
               openright,               					% inizio capitoli a destra
               english,                 
               italian          
               ]{book}    

%**************************************************************
% Importazione package
%************************************************************** 

%\usepackage{amsmath,amssymb,amsthm}    					% matematica

\usepackage[T1]{fontenc}                					% codifica dei font:
                                        					% NOTA BENE! richiede una distribuzione *completa* di LaTeX

\usepackage[utf8]{inputenc}             					% codifica di input; anche [latin1] va bene
                                        					% NOTA BENE! va accordata con le preferenze dell'editor

\usepackage[english, italian]{babel}    					% per scrivere in italiano e in inglese;
                                        					% l'ultima lingua (l'italiano) risulta predefinita

\usepackage{bookmark}                   					% segnalibri

\usepackage{caption}                    					% didascalie

\usepackage{chngpage, calc}             					% centra il frontespizio

\usepackage{csquotes}                   					% gestisce automaticamente i caratteri (")

\usepackage{colortbl}										% colorazione tabelle

\usepackage{emptypage}                  					% pagine vuote senza testatina e piede di pagina

\usepackage{epigraph}										% per epigrafi

\usepackage{eurosym}                    					% simbolo dell'euro

\usepackage{float}											% posizionamento di immagini e tabelle

%\usepackage{indentfirst}               					% rientra il primo paragrafo di ogni sezione

\usepackage{graphicx}                   					% immagini

\usepackage{hyperref}                  						% collegamenti ipertestuali

\usepackage[binding=5mm]{layaureo}      					% margini ottimizzati per l'A4; rilegatura di 5 mm

\usepackage{listings}                   					% codici

\usepackage{microtype}                  					% microtipografia

\usepackage{mparhack, fixltx2e, relsize}					% finezze tipografiche

\usepackage{nameref}                    					% visualizza nome dei riferimenti                                      

\usepackage[font=small]{quoting}        					% citazioni

\usepackage{subfig}                     					% sottofigure, sottotabelle

\usepackage[italian]{varioref}          					% riferimenti completi della pagina

\usepackage[[table, x11names, dvipsnames, table]{xcolor}	% colori

\usepackage{booktabs}                   					% tabelle                                       
\usepackage{tabularx}                   					% tabelle di larghezza prefissata                                    
\usepackage{longtable}                  					% tabelle su più pagine                                        
\usepackage{ltxtable}                   					% tabelle su più pagine e adattabili in larghezza

\usepackage[T1]{fontenc}
\usepackage{lmodern}

\usepackage[toc, acronym]{glossaries}   					% glossario
                                        					% per includerlo nel documento bisogna:
                                        					% 1. compilare una prima volta tesi.tex;
                                        					% 2. eseguire: makeindex -s tesi.ist -t tesi.glg -o tesi.gls tesi.glo
                                        					% 3. eseguire: makeindex -s tesi.ist -t tesi.alg -o tesi.acr tesi.acn
                                        					% 4. compilare due volte tesi.tex.

\usepackage[backend=biber, style=verbose-ibid, hyperref, backref]{biblatex}
                                        					% eccellente pacchetto per la bibliografia; 
                                       						% produce uno stile di citazione autore-anno; 
                                        					% lo stile "numeric-comp" produce riferimenti numerici
                                        					% per includerlo nel documento bisogna:
                                        					% 1. compilare una prima volta tesi.tex;
                                        					% 2. eseguire: biber tesi
                                        					% 3. compilare ancora tesi.tex.

%**************************************************************
% file contenente le impostazioni della tesi
%**************************************************************

%**************************************************************
% Frontespizio
%**************************************************************

% Autore
\newcommand{\myName}{Lorenzo Dei Negri}                                    
\newcommand{\myTitle}{Identità digitali per scambio e firma di documenti fiscali tramite blockchain}

% Tipo di tesi                   
\newcommand{\myDegree}{Tesi di laurea triennale}

% Università             
\newcommand{\myUni}{Università degli Studi di Padova}

% Facoltà       
\newcommand{\myFaculty}{Corso di Laurea in Informatica}

% Dipartimento
\newcommand{\myDepartment}{Dipartimento di Matematica "Tullio Levi-Civita"}

% Titolo del relatore
\newcommand{\profTitle}{Prof.}

% Relatore
\newcommand{\myProf}{Massimo Marchiori}

% Luogo
\newcommand{\myLocation}{Padova}

% Anno accademico
\newcommand{\myAA}{2019-2020}

% Data discussione
\newcommand{\myTime}{23 Luglio 2020}


%**************************************************************
% Impostazioni di impaginazione
% see: http://wwwcdf.pd.infn.it/AppuntiLinux/a2547.htm
%**************************************************************

\setlength{\parindent}{14pt}   % larghezza rientro della prima riga
\setlength{\parskip}{0pt}   % distanza tra i paragrafi


%**************************************************************
% Impostazioni di biblatex
%**************************************************************
\bibliography{bibliografia} % database di biblatex 

\defbibheading{bibliography} {
    \cleardoublepage
    \phantomsection 
    \addcontentsline{toc}{chapter}{\bibname}
    \chapter*{\bibname\markboth{\bibname}{\bibname}}
}

\setlength\bibitemsep{1.5\itemsep} % spazio tra entry

\DeclareBibliographyCategory{opere}
\DeclareBibliographyCategory{web}

\addtocategory{opere}{womak:lean-thinking}
\addtocategory{web}{site:agile-manifesto}

\defbibheading{opere}{\section*{Riferimenti bibliografici}}
\defbibheading{web}{\section*{Siti Web consultati}}


%**************************************************************
% Impostazioni di caption
%**************************************************************
\captionsetup{
    tableposition=top,
    figureposition=bottom,
    font=small,
    format=hang,
    labelfont=bf
}

%**************************************************************
% Impostazioni di glossaries
%**************************************************************
%**************************************************************
% Acronimi
%**************************************************************
\renewcommand{\acronymname}{Acronimi e abbreviazioni}

\newacronym[description=Ahead Of Time]
	{aot}{AOT}{Ahead Of Time}

\newacronym[description={\glslink{apig}{Application Program Interface}}]
    {api}{API}{Application Program Interface}
    
\newacronym[description=Commercio Cash Credit]
	{ccc}{CCC}{Commercio Cash Credit}
      
\newacronym[description=Commercio Token]
	{com}{COM}{Commercio Token}
    
\newacronym[description={\glslink{ddog}{DID Document}}]
	{ddo}{DDO}{DID Document}
    
\newacronym[description={\glslink{didg}{Decentralized Identifier}}]
    {did}{DID}{Decentralized Identifier}
    
\newacronym[description={\glslink{ictg}{Information and Communication Technologies}}]
	{ict}{ICT}{Information and Communication Technologies}
	
\newacronym[description=Just In Time]
	{jit}{JIT}{Just In Time}
	
\newacronym[description=Portable Document Format]
	{pdf}{PDF}{Portable Document Format}	

\newacronym[description={\glslink{pecg}{Posta Elettronica Certificata}}]
	{pec}{PEC}{Posta Elettronica Certificata}
	
\newacronym[description={\glslink{pocg}{Proof Of Concept}}]
	{poc}{POC}{Proof Of Concept}
	
\newacronym[description=Proof Of Stake]
	{pos}{POS}{Proof Of Stake}
	
\newacronym[description={\glslink{restg}{Representational State Transfer}}]
	{rest}{REST}{Representational State Transfer}
	
\newacronym[description={\glslink{sdkg}{Software Development Kit}}]
	{sdk}{SDK}{Software Development Kit}
	
\newacronym[description=Short Message Service]
	{sms}{SMS}{Short Message Service}

\newacronym[description={\glslink{spidg}{Sistema Pubblico d'Identità Digitale}}]
	{spid}{SPID}{Sistema Pubblico d'Identità Digitale}
	
\newacronym[description={\glslink{ssig}{Self-Sovereign Identity}}]
	{ssi}{SSI}{Self-Sovereign Identity}

\newacronym[description={\glslink{umlg}{Unified Modeling Language}}]
    {uml}{UML}{Unified Modeling Language}

%**************************************************************
% Glossario
%**************************************************************
\renewcommand{\glossaryname}{Glossario}

\newglossaryentry{apig}
{
    name=\glslink{api}{API},
    text=Application Program Interface,
    sort=api,
    description={in informatica con il termine \emph{Application Programming Interface (API)} (ing. interfaccia di programmazione di un'applicazione) si indica ogni insieme di procedure disponibili al programmatore, di solito raggruppate a formare una collezione di strumenti specifici per l'espletamento di un determinato compito all'interno di un certo programma. La finalità è ottenere un'astrazione, di solito tra l'hardware e il programmatore o tra software a basso e quello ad alto livello semplificando così il lavoro di programmazione}
}

\newglossaryentry{blockchaing}
{
	name=Blockchain,
	text=blockchain,
	sort=blockchain,
	description={in informatica con il termine \emph{blockchain} (ing. catena di blocchi) si indica una struttura dati condivisa e immutabile, le cui informazioni sono raggruppate in blocchi concatenati l'uno all'altro, a partire dal blocco d'origine, in ordine crescente per data di creazione, e la cui integrità viene garantita mediante l'uso della crittografia. Sebbene il numero di blocchi che la compongono sia destinato a crescere nel tempo, i dati contenuti al suo interno sono immutabili in quanto, una volta scritti non è più possibile né modificarli né eliminarli, a meno di non invalidare l'intera struttura. Il registro distribuito su cui la struttura si basa può essere acceduto, sia in lettura che in scrittura, da parte di tutti i nodi presenti nella rete che utilizza la \emph{blockchain} e, per garantire la coerenza tra le varie copie locali, l'aggiunta di un nuovo blocco viene gestita tramite un apposito protocollo di consenso condiviso}
}

\newglossaryentry{ddog}
{
	name=\glslink{ddo}{DDO},
	text=\glslink{did}{DID} Document,
	sort=ddo,
	description={in informatica con il termine \emph{\glslink{did}{DID} Document (DDO)} (ing. documento di un \glslink{did}{DID}) si indica un insieme di informazioni che descrivono il soggetto identificato da un \glslink{did}{DID}, corredati da opportuni meccanismi, come chiavi crittografiche e dati biometrici, attraverso i quali il soggetto può autenticare se stesso e dare prova della sua associazione al \glslink{did}{DID}}
}

\newglossaryentry{didg}
{
	name=\glslink{did}{DID},
	text=Decentralized Identifier,
	sort=did,
	description={in informatica con il termine \emph{Decentralized Identifier (DID)} (ing. identificativo decentralizzato) si indica un nuovo tipo di identificatore globale univoco che permette la creazione di identità digitali decentralizzate verificabili. Un \emph{DID} può identificare qualsiasi soggetto che necessiti di possedere un identificativo digitale e, a differenza delle altre identità digitali, non richiede la certificazione da parte di autorità centralizzate terze, fornitrici di attestazioni d'identità, per essere emesso; la registrazione tramite tecnologie \gls{blockchaing} o altre forme di reti decentralizzate, infatti, permette di risalire alle informazioni del soggetto identificato, senza bisogno di un'ente o di un'entità garante della loro veridicita e affidabilità}
}

\newglossaryentry{frameworkg}
{
	name=Framework,
	text=framework,
	sort=framework,
	description={in informatica con il termine \emph{framework} (ing. struttura o intelaiatura) si indica un'astrazione \textit{software} che fornisce diverse funzionalità generiche per lo sviluppo e la messa in produzione di applicazioni, permettendo comunque al programmatore di poterle modificare in base alle specifiche necessità del progetto. Rappresenta un ambiente \textit{software} universale e riutilizzabile che offre una piattaforma per facilitare lo sviluppo di applicazioni. Oltre a questo, può contenere al suo interno programmi di supporto, compilatori, librerire e \glslink{api}{API}}
}

\newglossaryentry{ictg}
{
	name=\glslink{ict}{ICT},
	text=Information and Communication Technologies,
	sort=ict,
	description={in informatica con il termine \emph{Information and Communication Technologies (ICT)} (ing. tecnologie di informazione e comunicazione) si indica ogni insieme di tecnologie riguardanti i sistemi integrati di telecomunicazione (linee di comunicazione cablate e senza fili), i \textit{computer}, le tecnologie audio-video e relativi \textit{software}, che permettono agli utenti di creare, immagazzinare e scambiare informazioni}
}

\newglossaryentry{namespaceg}
{
	name=Namespace,
	text=namespace,
	sort=namespace,
	description={in informatica con il termine \emph{namespace} (ing. spazio dei nomi) si indica un insieme di simboli, i nomi, che vengono utilizzati per identificare e riferire degli oggetti di vario genere. Lo scopo principale di un \emph{namespace} è assicurare che tutte le entità contenute al suo interno abbiano dei nomi univoci, in modo da poter essere facilmente identificabili. Gli spazi dei nomi sono spesso strutturati in gerarchie, in modo da poter riusare lo stesso nome per identificare oggetti distinti in contesti diversi.\\
	Nei linguaggi di programmazione, i \emph{namespace} sono tipicamente utilizzati per raggruppare un insieme di simboli e identificatori comuni ad una specifica funzionalità, in modo da evitare collisioni, ossia nomi duplicati che entrano in conflitto tra loro}
}

\newglossaryentry{nodo-validatoreg}
{
	name=Nodo validatore,
	text=nodo validatore,
	sort=nodo validatore,
	description={con il termine \emph{nodo validatore} si indica un nodo della rete Commercio.network che, oltre a mantenere in memoria tutto lo storico del \textit{database} distribuito che compone la \glslink{blockchaing}{blockchain}, ha anche la capacità di validare nuove transazioni che dovranno poi essere aggiunte in un nuovo blocco della catena. La validazione avviene tramite l'apposizione di una firma fatta con la propria chiave crittografica privata}
}

\newglossaryentry{open-sourceg}
{
	name=Open-source,
	text=open-source,
	sort=open-source,
	description={in informatica con il termine \emph{open-source} (ing. sorgente libero) si indica un \textit{software} di cui l'utente finale ha libero accesso a tutto il codice sorgente; questo può essere anche modificato per cambiarne e/o evolverne il funzionamento oppure per correggerne eventuali errori riscontrati durante l'utilizzo. L'utente inoltre ha anche la possibilità di ridistribuire pubblicamente la versione da lui rielaborata, infatti, rilasciando un \textit{software} con una licenza \emph{open-source}, gli sviluppatori rinunciano a qualsiasi diritto di proprietà intelluttuale sullo stesso}
}

\newglossaryentry{pecg}
{
	name=\glslink{pec}{PEC},
	text=Posta Elettronica Certificata,
	sort=pec,
	description={in informatica con il termine \emph{Posta Elettronica Certificata (PEC)} si indica il sistema di posta elettronica che fornisce al mittente ricevuta con valore legale dell’invio e della ricezione di messaggi informatici. La \textit{PEC} è quindi un sistema in grado di certificare l'identità (digitale) di un soggetto che ne fa utilizzo, garantendo inoltre il non ripudio per tutte le operazioni da esso svolte}
}

\newglossaryentry{pocg}
{
	name=\glslink{poc}{POC},
	text=proof of concept,
	sort=poc,
	description={in informatica con il termine \emph{Proof Of Concept (POC)} (ing. prova di concetto) si indica la realizzazione di un \textit{software} allo scopo di dimostrare la fattibilità di un certo principio o di una certa idea che sta alla base di un progetto \textit{software} più ampio e completo. Un \emph{Proof Of Concept (POC)} è solitamente di dimensioni ridotte e può anche essere incompleto, questo perché il suo fine è solamente verificare il potenziale pratico di un concetto teorico}
}

\newglossaryentry{power-upg}
{
	name=Powerup,
	text=powerup,
	sort=powerup,
	description={nel contesto della Commercio.network, con il termine \emph{powerup} (ing. potenziamento), si indica l'operazione attraverso cui un nodo della rete, identificato da un proprio \glslink{did}{DID}, richiede il trasferimento di un certo ammontare di \glslink{tokeng}{token} dal proprio \glslink{walletg}{wallet} verso quello di un altro nodo, identificato a sua volta da un suo \glslink{did}{DID}}
}

\newglossaryentry{restg}
{
	name=\glslink{rest}{REST},
	text=Representational State Transfer,
	sort=rest,
	description={in informatica con il termine \emph{Representational State Transfer (REST)} (ing. trasferimento di stato rappresentativo) si indica uno stile architetturale \textit{software} che definisce un insieme di vincoli da rispettare nella creazione di servizi \textit{web}. I servizi \textit{web} conformi a questi principi vengono chiamati \textit{RESTful} e forniscono un alto grado di interoperabilità tra i diversi sistemi che comunicano in rete. Infatti permettono di accedere e di elaborare una rappresentazione testuale delle risorse esposte dai servizi stessi, utilizzando un insieme di operazioni uniformi e predefinite.\\
	Ogni risorsa, identificata da uno specifico nome univoco, può essere acceduta sia in lettura che in scrittura, in base ai permessi previsti per la stessa, attraverso appositi metodi del protocollo utilizzato per la gestione delle richieste da parte dei \textit{client} e le risposte da parte del \textit{server}}
}

\newglossaryentry{sdkg}
{
	name=\glslink{sdk}{SDK},
	text=Software Development Kit,
	sort=sdk,
	description={in informatica con il termine \emph{Software Development Kit (SDK)} (ing. pacchetto di sviluppo per \textit{software}) si indica una collezione di strumenti per lo sviluppo \textit{software} contenuti all'interno di un singolo pacchetto installabile all'interno del proprio sistema; questi permettono di facilitare, a volte notevolmente, la creazione di applicazioni. Questi strumenti solitamente sono specifici per il particolare tipo di \textit{hardware}, sistema operativo e linguaggio di programmazione utilizzati per lo sviluppo \textit{software}}
}

\newglossaryentry{spidg}
{
	name=\glslink{spid}{SPID},
	text=Sistema Pubblico d'Identità Digitale,
	sort=spid,
	description={con il termine \emph{Sistema Pubblico d'Identità Digitale (SPID)} si indica la piattaforma tramite cui le pubbliche amministrazioni consentono ai cittadini e alle imprese l'accesso in rete ai propri servizi tramite un unico sistema di \textit{password}. Questo sistema consente quindi di identitificare univocamente ogni soggetto che ne fa utilizzo, fornendogli un'identità digitale che lo rappresenta e che gli permette di usufruire di tutti quei servizi che richiedono un'autenticazione tramite documento d'identificazione}
}

\newglossaryentry{ssig}
{
	name=\glslink{ssi}{SSI},
	text=Self-Sovereign Identity,
	sort=ssi,
	description={in informatica con il termine \emph{Self-Sovereign Identity (SSI)} (ing. identità sovrana) si indica un meccanismo che permette agli individui e agli enti di creare la propria identità digitale, mantenendo il pieno controllo delle realtive credenziali, senza bisogno di richiederne l'emissione certificata da parte di un intermediaro terzo o di un'autorità centralizzata. I titolari delle identità hanno la totale padronanza dei propri dati personali e possono decidere in autonomia in che termini questi possono essere utilizzati e condivisi}
}

\newglossaryentry{tokeng}
{
	name=Token,
	text=token,
	sort=token,
	description={in informatica, all'interno del contesto di una \glslink{blockchaing}{blockchain}, con il termine \emph{token} (ing. gettone) si indica un bene digitale che può essere visto come una collezione di informazioni in grado di garantire un diritto di proprietà, al soggetto possessore, su un certo insieme di dati presenti all'interno della \glslink{blockchaing}{blockchain} stessa. I \emph{token} possono essere scambiati tra due diversi partecipanti alla rete senza bisogno di un intermediario, ma semplicemente effettuando una transazione; infatti questi non sono emessi da un ente, ma fanno nativamente parte della \glslink{blockchaing}{blockchain} in cui vengono generati.\\
	Oltre a questo, possono essere utilizzati come delle vere e proprie monete digitali: in questo caso assumono il nome di criptovalute in quanto la loro progettazione si basa in gran parte su tecnologie crittografiche. Il loro valore viene dettato dalle leggi economiche del mercato, in particolare da quella della domanda e dell'offerta, applicata ai nodi della rete; infatti il loro controvalore in valuta legale è spesso molto variabile}
}

\newglossaryentry{umlg}
{
    name=\glslink{uml}{UML},
    text=UML,
    sort=uml,
    description={in ingegneria del \textit{software} l'\emph{Unified Modeling Language (UML)} (ing. linguaggio di modellazione unificato) è un linguaggio di modellazione e specifica basato sul paradigma \textit{object-oriented}. L'\emph{UML} svolge un'importantissima funzione di ``lingua franca'' nella comunità della progettazione e programmazione a oggetti. Gran parte della letteratura di settore usa tale linguaggio per descrivere soluzioni analitiche e progettuali in modo sintetico e comprensibile a un vasto pubblico}
}

\newglossaryentry{walletg}
{
	name=Wallet,
	text=wallet,
	sort=wallet,
	description={in informatica, all'interno del contesto di una \glslink{blockchaing}{blockchain}, con il termine \emph{wallet} (ing. portafoglio) si indica un sistema \textit{hardware} o \textit{software} che permette di tracciare l'ammontare del numero di \glslink{tokeng}{token}, o criptovalute, possedute da un utente della \glslink{blockchaing}{blockchain}. In realtà il \emph{wallet} non contiene effettivamente al suo interno le criptovalute, queste sono invece memorizzate dentro al \textit{database} distribuito che costituisce la \glslink{blockchaing}{blockchain} stessa; infatti, attraverso il \emph{wallet} è possibile solamente scambiare (inviare o ricevere) dei \glslink{tokeng}{token}, con gli altri utenti partecipanti alla rete, tramite delle transazioni. Per fare questo viene mantenuta una coppia di chiavi crittografiche, una privata ed una pubblica, all'interno del \emph{wallet} stesso: la chiave pubblica serve per ricevere dei versamenti, mentre la chiave privata viene utilizzata per effettuare dei pagamenti dal proprio \emph{wallet}}
}
 % database di termini
\makeglossaries


%**************************************************************
% Impostazioni di graphicx
%**************************************************************
\graphicspath{{immagini/}} % cartella dove sono riposte le immagini


%**************************************************************
% Impostazioni di hyperref
%**************************************************************
\hypersetup{
    %hyperfootnotes=false,
    %pdfpagelabels,
    %draft,	% = elimina tutti i link (utile per stampe in bianco e nero)
    colorlinks=true,
    linktocpage=true,
    pdfstartpage=1,
    pdfstartview=FitV,
    % decommenta la riga seguente per avere link in nero (per esempio per la stampa in bianco e nero)
    %colorlinks=false, linktocpage=false, pdfborder={0 0 0}, pdfstartpage=1, pdfstartview=FitV,
    breaklinks=true,
    pdfpagemode=UseNone,
    pageanchor=true,
    pdfpagemode=UseOutlines,
    plainpages=false,
    bookmarksnumbered,
    bookmarksopen=true,
    bookmarksopenlevel=1,
    hypertexnames=true,
    pdfhighlight=/O,
    %nesting=true,
    %frenchlinks,
    urlcolor=webbrown,
    linkcolor=RoyalBlue,
    citecolor=webgreen,
    %pagecolor=RoyalBlue,
    %urlcolor=Black, linkcolor=Black, citecolor=Black, %pagecolor=Black,
    pdftitle={\myTitle},
    pdfauthor={\textcopyright\ \myName, \myUni, \myFaculty},
    pdfsubject={},
    pdfkeywords={},
    pdfcreator={pdfLaTeX},
    pdfproducer={LaTeX}
}

%**************************************************************
% Impostazioni di itemize
%**************************************************************
\renewcommand{\labelitemi}{$\ast$}

%\renewcommand{\labelitemi}{$\bullet$}
%\renewcommand{\labelitemii}{$\cdot$}
%\renewcommand{\labelitemiii}{$\diamond$}
%\renewcommand{\labelitemiv}{$\ast$}


%**************************************************************
% Impostazioni di listings
%**************************************************************
\lstset{
    language=[LaTeX]Tex,%C++,
    keywordstyle=\color{RoyalBlue}, %\bfseries,
    basicstyle=\small\ttfamily,
    %identifierstyle=\color{NavyBlue},
    commentstyle=\color{Green}\ttfamily,
    stringstyle=\rmfamily,
    numbers=none, %left,%
    numberstyle=\scriptsize, %\tiny
    stepnumber=5,
    numbersep=8pt,
    showstringspaces=false,
    breaklines=true,
    frameround=ftff,
    frame=single
} 


%**************************************************************
% Impostazioni di xcolor
%**************************************************************
\definecolor{webgreen}{rgb}{0,.5,0}
\definecolor{webbrown}{rgb}{.6,0,0}


%**************************************************************
% Altro
%**************************************************************

\newcommand{\omissis}{[\dots\negthinspace]} % produce [...]

% eccezioni all'algoritmo di sillabazione
\hyphenation
{
    ma-cro-istru-zio-ne
    gi-ral-din
}

\newcommand{\sectionname}{sezione}
\addto\captionsitalian{\renewcommand{\figurename}{Figura}
                       \renewcommand{\tablename}{Tabella}}

\newcommand{\glsfirstoccur}{\ap{{[g]}}}

\newcommand{\intro}[1]{\emph{\textsf{#1}}}

%**************************************************************
% Environment per ``rischi''
%**************************************************************
\newcounter{riskcounter}                % define a counter
\setcounter{riskcounter}{0}             % set the counter to some initial value

%%%% Parameters
% #1: Title
\newenvironment{risk}[1]{
    \refstepcounter{riskcounter}        % increment counter
    \par \noindent                      % start new paragraph
    \textbf{\arabic{riskcounter}. #1}   % display the title before the 
                                        % content of the environment is displayed 
}{
    \par\medskip
}

\newcommand{\riskname}{Rischio}

\newcommand{\riskdescription}[1]{\textbf{\\Descrizione:} #1.}

\newcommand{\risksolution}[1]{\textbf{\\Soluzione:} #1.}

%**************************************************************
% Environment per ``use case''
%**************************************************************
\newcounter{usecasecounter}             % define a counter
\setcounter{usecasecounter}{0}          % set the counter to some initial value

%%%% Parameters
% #1: ID
% #2: Nome
\newenvironment{usecase}[2]{
    \renewcommand{\theusecasecounter}{\usecasename #1}  % this is where the display of 
                                                        % the counter is overwritten/modified
    \refstepcounter{usecasecounter}             % increment counter
    \vspace{10pt}
    \par \noindent                              % start new paragraph
    {\large \textbf{\usecasename #1: #2}}       % display the title before the 
                                                % content of the environment is displayed 
    \medskip
}{
    \medskip
}

\newcommand{\usecasename}{UC}

\newcommand{\usecaseactors}[1]{\textbf{\\Attori Principali:} #1. \vspace{4pt}}
\newcommand{\usecasepre}[1]{\textbf{\\Precondizioni:} #1. \vspace{4pt}}
\newcommand{\usecasedesc}[1]{\textbf{\\Descrizione:} #1. \vspace{4pt}}
\newcommand{\usecasepost}[1]{\textbf{\\Postcondizioni:} #1. \vspace{4pt}}
\newcommand{\usecasealt}[1]{\textbf{\\Scenario Alternativo:} #1. \vspace{4pt}}

%**************************************************************
% Environment per ``namespace description''
%**************************************************************

\newenvironment{namespacedesc}{
    \vspace{10pt}
    \par \noindent                              % start new paragraph
    \begin{description} 
}{
    \end{description}
    \medskip
}

\newcommand{\classdesc}[2]{\item[\textbf{#1:}] #2}                     					% file con le impostazioni personali

\begin{document}
	
%**************************************************************
% Materiale iniziale
%**************************************************************

\frontmatter
% !TEX encoding = UTF-8
% !TEX TS-program = pdflatex
% !TEX root = ../tesi.tex

%**************************************************************
% Frontespizio 
%**************************************************************
\begin{titlepage}

\begin{center}

\begin{LARGE}
\textbf{\myUni}\\
\end{LARGE}

\vspace{10pt}

\begin{Large}
\textsc{\myDepartment}\\
\end{Large}

\vspace{10pt}

\begin{large}
\textsc{\myFaculty}\\
\end{large}

\vspace{30pt}
\begin{figure}[htbp]
\begin{center}
\includegraphics[height=6cm]{logo-unipd}
\end{center}
\end{figure}
\vspace{30pt} 

\begin{LARGE}
\begin{center}
\textbf{\myTitle}\\
\end{center}
\end{LARGE}

\vspace{10pt} 

\begin{large}
\textsl{\myDegree}\\
\end{large}

\vspace{40pt} 

\begin{large}
\begin{flushleft}
\textit{Relatore}\\ 
\vspace{5pt} 
\profTitle{} \myProf
\end{flushleft}

\vspace{0pt} 

\begin{flushright}
\textit{Laureando}\\ 
\vspace{5pt} 
\myName
\end{flushright}
\end{large}

\vspace{40pt}

\line(1, 0){338} \\
\begin{normalsize}
\textsc{Anno Accademico \myAA}
\end{normalsize}

\end{center}
\end{titlepage} 
\input{inizio-fine/colophon}
% !TEX encoding = UTF-8
% !TEX TS-program = pdflatex
% !TEX root = ../tesi.tex

%**************************************************************
% Dedica
%**************************************************************
\cleardoublepage
\phantomsection
\thispagestyle{empty}
\pdfbookmark{Dedica}{Dedica}

\vspace*{3cm}

\begin{center}
\textit{"Man is still the most extraordinary computer of all."} \\ \medskip
--- John Fitzgerald Kennedy
\end{center}

\medskip

\begin{center}
Dedicato alla mia famiglia, a Fiammetta e a tutti i miei amici.
\end{center}

% !TEX encoding = UTF-8
% !TEX TS-program = pdflatex
% !TEX root = ../tesi.tex

%**************************************************************
% Sommario
%**************************************************************
\cleardoublepage
\phantomsection
\pdfbookmark{Sommario}{Sommario}
\begingroup
\let\clearpage\relax
\let\cleardoublepage\relax
\let\cleardoublepage\relax

\chapter*{Sommario}

Il presente documento descrive il lavoro svolto durante il periodo di stage, della durata di trecento ore, dal laureando \myName{} presso l'azienda \myCompany{} \companyTitle{}, situata a Limena (PD).\\
Lo stage ha avuto come argomento principale la \textit{blockchain}, le tecnologie ad essa correlate e come queste possono essere utilizzate per la creazione di identità digitali e la condivisione di documenti firmati. Il progetto, nella sua interezza, è di particolare rilevanza per l'azienda, in quanto direttamente coinvolto nel raggiungimento della sua \textit{mission}: eliminare in maniera definitiva la generazione di documenti cartacei all'interno delle aziende.
Gli obiettivi da raggiungere erano molteplici e sono stati divisi in due parti distinte, una teorica di ricerca e formazione sull'argomento trattato e una pratica di progettazione e sviluppo di un relativo \textit{proof of concept}.\\
Nel primo periodo, dopo la necessaria formazione sul contesto aziendale, era richiesto lo studio teorico di tutti i concetti legati alla \textit{blockchain}, analizzandone il funzionamento, lo scopo, i possibili utilizzi in casi reali, i pregi e i difetti; successivamente bisognava analizzare, allo stesso modo, la \textit{blockchain} sviluppata da Commerc.io S.r.l., una start-up innovativa associata a \myCompany{} \companyTitle{} che punta alla realizzazione di una rete di aziende che consenta la firma e lo scambio di documenti tra i partecipanti. Infine, in preparazione alla realizzazione del \textit{proof of concept} previsto dal progetto, era necessario approfondire la documentazione della \textit{blockchain} \textit{Commercio.network} dal punto di vista degli strumenti di sviluppo, ossia quali librerie sono disponibili per interagire con la rete stessa.\\
Nel secondo periodo era richiesta l'analisi, la progettazione e l'implementazione di un'applicazione \textit{mobile}, scritta in linguaggio \textit{Dart} attraverso il \textit{framework} \textit{Flutter}, per la creazione di identità digitali e la condivisione di documenti tramite la rete \textit{Commercio.network}. Inoltre, per raggiungere tale obiettivo, è stato necessario un periodo di formazione individuale sugli strumenti di sviluppo scelti dall'azienda per la realizzazione dell'applicazione.

%\vfill
%
%\selectlanguage{english}
%\pdfbookmark{Abstract}{Abstract}
%\chapter*{Abstract}
%
%\selectlanguage{italian}

\endgroup			

\vfill


% !TEX encoding = UTF-8
% !TEX TS-program = pdflatex
% !TEX root = ../tesi.tex

%**************************************************************
% Ringraziamenti
%**************************************************************
\cleardoublepage
\phantomsection
\pdfbookmark{Ringraziamenti}{ringraziamenti}

\begin{flushright}{
	\slshape    
	``Anyone who has lost track of time when using a computer knows the propensity to dream, the urge to make dreams come true and the tendency to miss lunch.''} \\ 
	\medskip
	--- Tim Berners-Lee
\end{flushright}

\bigskip

\begingroup
\let\clearpage\relax
\let\cleardoublepage\relax
\let\cleardoublepage\relax

\chapter*{Ringraziamenti}

\noindent \textit{Innanzitutto, vorrei esprimere la mia gratitudine al \profTitle{} \myProf, relatore della mia tesi, sia per l'aiuto e il sostegno fornitomi durante la stesura del lavoro e nel corso del progetto di stage, sia per la passione trasmessa nello studio di alcuni temi dell'informatica. La dedizione nei confronti del suo lavoro è stata per me stimolo ed ispirazione.}\\

\noindent \textit{Desidero ringraziare con affetto la mia famiglia per il supporto, anche economico, e il grande aiuto che mi hanno dato durante tutto il periodo di studi. La stabilità, che non mi hanno mai fatto mancare, mi ha permesso di affrontare il percorso universitario con la serenità e la concentrazione necessaria.}\\

\noindent \textit{Voglio dedicare un pensiero speciale a Fiammetta per lo splendido rapporto che abbiamo costruito durante questi anni e per essermi stata sempre vicina, sostenendomi nei momenti difficili e condividendo insieme quelli più felici.}\\

\noindent \textit{Infine, desidero ringraziare di cuore tutti i miei amici per i bellissimi momenti passati insieme durante questi tre anni e le mille avventure vissute. Un ringraziamento particolare va a Massimo per essere sempre stato pronto ad aiutarmi nel momento del bisogno, sia nella sfera universitaria che in quella personale, e per aver superato insieme ogni fatica ed ogni ostacolo che abbiamo trovato durante il nostro percorso di studi.}\\
\bigskip

\noindent\textit{\myLocation, \myTime}
\hfill \myName

\endgroup


\input{inizio-fine/indici}
\cleardoublepage

%**************************************************************
% Materiale principale
%**************************************************************

\mainmatter
% !TEX encoding = UTF-8
% !TEX TS-program = pdflatex
% !TEX root = ../tesi.tex

%**************************************************************
\chapter{Introduzione}
\label{cap:introduzione}

%**************************************************************
\section{L'azienda}

\myCompany{} \companyTitle{}\footcite{site:2c-solution} è un'azienda italiana con sede a Limena (PD) che opera nel settore della digitalizzazione aziendale da circa quindici anni. La sua \textit{mission} è far cessare definitivamente l'utilizzo di documenti cartacei all'interno delle aziende, in modo da semplificare le procedure documentali, agevolando e tracciando il trasferimento delle informazioni; tutto questo nel pieno rispetto delle normative vigenti in materia di \textit{privacy} e sicurezza informatica.

\begin{figure}[!h] 
	\centering 
	\includegraphics[width=0.75\columnwidth]{logo-2c-solution} 
	\caption{Logo di \myCompany{} \companyTitle}
\end{figure}

L'azienda nasce dall'idea di due professionisti dell'\gls{ict}\glsfirstoccur{} di creare un'unica soluzione per permettere alle imprese di eliminare tutta la carta dalla gestione aziendale. Successivamente il prodotto è stato ampliato con diversi moduli, ciascuno dedicato alla realizzazione di una delle tre funzionalità principali:

\begin{itemize}
	\item gestione dei processi aziendali ad alte prestazioni;
	\item organizzazione, catalogazione e ricerca dei documenti tramite una piattaforma digitale semplice ma potente;
	\item creazione di portali accessibili in rete per la condivisione dei documenti aziendali.
\end{itemize}

Col passare del tempo e con l'aumento del numero di regolamentazioni in merito alla gestione digitale delle imprese, l'azienda ha deciso di concentrarsi sulla realizzazione di un nuovo prodotto in grado di consentire la completa gestione documentale a norma di legge. Anche questo \textit{software} è stato diviso in più moduli, in quanto i servizi che offre sono molteplici, in particolare sono implementate le seguenti funzionalità:

\begin{itemize}
	\item firma digitale, archiviazione e conservazione di documenti con valenza legale;
	\item fatturazione elettronica europea, per professionisti, aziende e per la pubblica amministrazione;
	\item gestione di identità digitali personali tramite \gls{spid}\glsfirstoccur{} e \gls{pec}\glsfirstoccur.
\end{itemize}

Fin dall'inizio il lavoro dell'azienda si è basato sui principi che hanno ispirato i due soci fondatori ad iniziare la realizzazione del loro progetto:

\begin{itemize}
	\item l'innovazione tecnologica nella ricerca, sviluppo ed utilizzo di soluzioni che facciano uso di strumenti e tecnologie all'avanguardia;
	\item l'aggiornamento continuo dei prodotti sviluppati e la formazione del personale rispetto alle nuove certificazioni e normative vigenti riguardo la gestione della documentazione e dei processi aziendali;
	\item la grande passione per il lavoro svolto e per i frutti che questo porta, fondamentale per svolgerlo al meglio con gioia ed energia.
\end{itemize}

%**************************************************************
\section{Lo stage}

Il principale pilastro che sorregge il lavoro svolto da \myCompany{} \companyTitle{} è la constante ricerca di innovazione e di nuove tecnologie da utilizzare per migliorare i propri prodotti o per offrirne di nuovi ai propri \textit{partner} e clienti. A tal proposito, l'azienda punta molto sulla ricerca di nuove soluzioni ai problemi che affronta quotidianamente e da quest'anno partecipa all'evento StageIT, in convenzione con l'Università di Padova, per presentare progetti di stage a ragazzi universitari promettenti, con la possibilità di una futura assunzione.\\
Il percorso di stage offerto prevedeva l'inserimento del candidato nel reparto di ricerca e sviluppo, affiancato dal tutor aziendale \myTutor, incaricato di supervisionare e guidare lo stagista durante il lavoro.\\
\myCompany{} \companyTitle{} punta ad inserire la tecnologia \gls{blockchaing}\glsfirstoccur{} nei propri sistemi di gestione dei documenti e delle identità digitali in modo da sfruttarne le proprietà per migliorare le funzionalità già offerte attualmente. A tale scopo l'azienda ha ritenuto che la soluzione migliore fosse l'utilizzo della \gls{blockchaing} realizzata da Commerc.io S.r.l.\footcite{site:commerc-io} (una società affiliata a \myCompany{} \companyTitle), la \textit{commercio.network}\footcite{site:commercio-network}: infatti questa rete \gls{open-sourceg}\glsfirstoccur{} è stata progettata appositamente per la condivisione di documenti tra le aziende che vi partecipano come nodi. Tuttavia questa non è l'unica funzionalità disponibile, infatti all'interno della \gls{blockchaing} è anche possibile gestire delle identità e delle firme digitali, tutto nel pieno rispetto delle normative europee vigenti in materia.\\
L'obiettivo è quindi realizzare un'applicazione che permetta agli utenti di creare e gestire una propria identità digitale in maniera completamente autonoma e, successivamente, che consenta loro il caricamento e la condivisione di documenti. Tuttavia, prima di procedere al suo sviluppo è necessario analizzare il funzionamento della \textit{commercio.network} per comprenderne pregi e difetti, oltre a capirne possibili le modalità d'interazione: quali linguaggi di programmazione possono essere utilizzati, quali sono le \gls{api}\glsfirstoccur{} esposte e quali \textit{software} e/o servizi devono essere utilizzati nell'ambiente d'esecuzione.\\
Di conseguenza, tenendo in considerazione che lo sviluppo della \textit{commercio.network} non è ancora completo, il progetto di stage proposto da \myCompany{} \companyTitle{} si articola in due parti distinte ma allo stesso tempo complementari:

\begin{itemize}
	\item primo periodo: dopo la preliminare formazione sul contesto aziendale, è richiesto lo studio delle tecnologie \gls{blockchaing} da un punto di vista generale, in preparazione all'analisi della rete \textit{commercio.network} per capirne le caratteristiche e valutarne l'impiego nel proseguo del progetto aziendale;
	\item secondo periodo: basandosi su quanto appreso nella precedente fase di studio, è richiesto lo sviluppo di un'applicazione con funzione di \gls{poc}\glsfirstoccur{}, per dimostrare la fattibilità del \textit{software} di gestione delle identità digitali e dei documenti associati, utilizzando tutte le tecnologie e gli strumenti necessari per accedere alle funzionalità della \textit{commercio.network}. 
 \end{itemize}

%**************************************************************
\section{Organizzazione del testo}

\begin{description}
    \item[{\hyperref[cap:stage]{Il secondo capitolo}}] descrive ...
    
    \item[{\hyperref[cap:blockchain]{Il terzo capitolo}}] approfondisce ...
    
    \item[{\hyperref[cap:tecnologie-strumenti]{Il quarto capitolo}}] approfondisce ...
    
    \item[{\hyperref[cap:applicazione]{Il quinto capitolo}}] approfondisce ...
    
    \item[{\hyperref[cap:conclusioni]{Il sesto capitolo}}] approfondisce ...
\end{description}

Per la stesura del documento sono state adottate le seguenti convenzioni tipografiche:

\begin{itemize}
	\item gli acronimi, le abbreviazioni e i termini ambigui o di uso non comune menzionati vengono definiti nel glossario, situato alla fine del presente documento;
	\item per la prima occorrenza dei termini riportati nel glossario viene utilizzata la seguente nomenclatura: \emph{termine}\glsfirstoccur;
	\item i termini in lingua straniera o facenti parti del gergo tecnico sono evidenziati con il carattere \emph{corsivo}.
\end{itemize}
             					% Introduzione
% !TEX encoding = UTF-8
% !TEX TS-program = pdflatex
% !TEX root = ../tesi.tex

%**************************************************************
\chapter{Stage}
\label{cap:processi-metodologie}
%**************************************************************

%\intro{Brevissima introduzione al capitolo}\\

\section{Progetto}

\myCompany{} \companyTitle{} fornisce diversi strumenti per eliminare l'utilizzo di documenti cartacei all’interno delle aziende, in modo da agevolare e tracciare il trasferimento delle informazioni. Oltre a \myCompany{} \companyTitle, il progetto coinvolge anche la start-up Commerc.io che gestisce lo scambio di documenti siglati da firme elettroniche tramite una propria blockchain.
\\
Lo scopo del progetto di stage è la realizzazione di un applicativo che veda l'integrazione, tramite la blockchain sviluppata da Commerc.io, di identità digitali basate su self-sovereign identity e sulla notarizzazione di processi di condivisione e sottoscrizione, con firme elettroniche, dei documenti fiscali.
\\
Con self-sovereign identity si intende un meccanismo che permette agli individui di certificare la propria identità digitale, senza bisogno che un garante esterno lo faccia per loro.
\\
Tutto questo dovrà essere integrato all'interno dei processi di digital transformation di \myCompany{} \companyTitle.
\\
Lo studente dovrà, inizialmente, effettuare uno studio ed un'analisi dettagliata della documentazione del progetto Commerc.io nel contesto dei documenti fiscali, in modo da approfondire le tecnologie riguardanti la blockchain. Dovrà poi essere sviluppato un \textit{proof of concept}, che si servirà delle API appositamente implementate da Commerc.io per l'interazione con la loro rete \textit{commercio.network}.
\\
L'attività di sviluppo prevederà, oltre all'analisi iniziale, un periodo di progettazione, che dovrà essere opportunamente argomentata e documentata. Il prototipo finale, risultante dall'attività di codifica, dovrà permettere l'interazione con la blockchain di Commerc.io per il riconoscimento delle identità digitali e la creazione di transazioni per la trasmissione di documenti.
\\
Durante questo progetto di stage lo studente avrà occasione di approfondire le sue conoscenze nell'ambito delle blockchain, le tecnologie che le riguardano ed il loro utilizzo in situazioni reali.
\\
In particolare lo studio si concentrerà sulla:
\begin{itemize}
	\item familiarizzazione col contesto aziendale di fatturazione elettronica;
	\item comprensione del funzionamento di una blockchain;
	\item integrazione di identità digitali tramite self-sovereign identity, sfruttando la blockchain;
	\item applicazione della blockchain per garantire l'integrità di documenti condivisi;
	\item applicazione della blockchain per implementare un servizio di decentralized identifier (DID).
\end{itemize}


\section{Obiettivi}

\subsection*{Classificazione}
Si farà riferimento ai requisiti secondo la seguente classificazione, la quale permette di identificarli univocamente:
\begin{itemize}
	\item \textbf{RO-X:} requisiti obbligatori, vincolanti in quanto obiettivo primario richiesto dall'azienda proponente;
	\item \textbf{RD-X:} requisiti desiderabili, non vincolanti o strettamente necessari, ma dal riconoscibile valore aggiunto;
	\item \textbf{RF-X:} requisiti facoltativi, rappresentanti valore aggiunto non strettamente competitivo.
\end{itemize}

Nelle sigle precedentemente indicate, \textbf{X} è un numero intero progressivo maggiore di zero con funzione di identificativo del requisito.

\subsection*{Definizione}
Si prevede lo raggiungimento dei seguenti obiettivi.

\subsubsection*{Obbligatori}
\begin{itemize}
	\item \textbf{RO-1:} acquisizione delle competenze sulle tematiche relative alla blockchain e alle tecnologie associate;
	\item \textbf{RO-2:} implementazione di un \textit{proof of concept} che permetta l'interazione con la blockchain di Commerc.io per il riconoscimento delle identità digitali e la trasmissione di documenti, tramite un'interfaccia grafica basilare;
	\item \textbf{RO-3:} stesura della documentazione relativa ai risultati ottenuti nel corso delle varie fasi del progetto e alla gestione delle difficoltà incontrate;
	\item \textbf{RO-4:} stesura della relazione riguardante il lavoro svolto durante tutto il periodo di stage.
\end{itemize}

\subsubsection*{Desiderabili}
\begin{itemize}
	\item \textbf{RD-1:} implementazione di un'interfaccia web completa per il \textit{proof of concept}.
\end{itemize}

\subsubsection*{Facoltativi}
\begin{itemize}
	\item \textbf{RF-1:} sviluppo di un'applicazione mobile in Flutter che offra le funzionalità richieste per il \textit{proof of concept}.
\end{itemize} 


\section{Prodotti}

Lo studente al termine del periodo si stage dovrà fornire i seguenti prodotti all'azienda proponente \myCompany{} \companyTitle.
\subsubsection*{Obbligatori}
\begin{itemize}
	\item il \textit{proof of concept} che mostri il funzionamento della gestione delle identità digitali e del trasferimento di documenti tramite \textit{blockchain}, che faccia uso delle API sviluppate da Commerc.io e offra una semplice interfaccia per il suo utilizzo;
	\item la documentazione dettagliata dei risultati ottenuti, comprendente un'analisi critica delle difficoltà riscontrate e delle soluzioni adottate;
	\item una relazione che illustri le attività svolte e i prodotti realizzati;
	\item una presentazione finale del lavoro svolto, da esporre in azienda.
\end{itemize}

\subsubsection*{Desiderabili}
\begin{itemize}
	\item un'interfaccia web completa che permetta l'utilizzo del \textit{proof of concept}.
\end{itemize}

\subsubsection*{Facoltativi}
\begin{itemize}
	\item un'applicazione mobile scritta in Flutter che offra le funzionalità richieste dal progetto.
\end{itemize} 


\section{Pianificazione}

Le attività di progetto saranno scandite dalla seguente pianificazione, per un totale di 300 ore.

\begin{tabularx}{\textwidth}{|c|c|c|p{7.85cm}|}
	\hline
	\textbf{Settimana} & \textbf{Date} & \textbf{Ore di lavoro} & \textbf{Descrizione dell'attività} \\\hline
	I & 06/05/2020 - 08/05/2020 & 20 & Studio dei concetti fondamentali legati alla blockchain e del contesto aziendale. \\\hline
	II & 11/05/2020 - 15/05/2020 & 40 & Studio approfondito della blockchain di Commerc.io e delle funzionalità da implementare. \\\hline
	III & 18/05/2020 - 22/05/2020 & 40 & Implementazione e verifica delle funzionalità offerte ad un account registrato. \\\hline
	IV & 25/05/2020 - 29/05/2020 & 40 & Implementazione e verifica delle funzionalità di gestione dell'identità digitale. \\\hline
	V & 01/06/2020 - 05/06/2020 & 40 & Implementazione e verifica delle funzionalità di condivisione di documenti. \\\hline
	VI & 08/06/2020 - 12/06/2020 & 40 & Implementazione dell'interfaccia grafica e validazione del prototipo. \\\hline
	VII & 15/06/2020 - 19/06/2020 & 40 & Stesura della documentazione ed inizio della relazione. \\\hline
	VIII & 22/06/2020 - 26/06/2020 & 40 & Termine stesura della relazione e presentazione di quanto svolto in azienda. \\\hline
\end{tabularx}

\subsection*{Prima settimana}
\begin{itemize}
	\item studio dei concetti fondamentali legati alla blockchain;
	\item formazione riguardo il contesto aziendale della fatturazione elettronica;
	\item studio della documentazione della blockchain sviluppata da Commerc.io.
\end{itemize}

\subsection*{Seconda settimana} 
\begin{itemize}
	\item studio della funzione di un account riferito alla blockchain di Commerc.io, con particolare attenzione al concetto di wallet;
	\item studio della creazione e del mantenimento di un'identità digitale all'interno della blockchain e delle sue possibili interazioni con essa;
	\item studio delle API sviluppate da Commerc.io per la gestione dei documenti tramite la blockchain.
\end{itemize}

\subsection*{Terza settimana}
\begin{itemize}
	\item implementazione delle seguenti funzionalità, relative ad un account nella blockchain di Commerc.io:
	\begin{itemize}
		\item generazione ed importazione di un HD Wallet;
		\item richiesta, ricezione ed invio di Token tra diversi account, all'interno della blockchain;
		\item controllo del bilancio di uno specifico account;
		\item acquisizione e perdita dello stato di Validator per i Token utilizzati nella blockchain;
	\end{itemize}
	\item verifica delle funzionalità implementate.
\end{itemize}

\subsection*{Quarta settimana} 
\begin{itemize}
	\item implementazione delle seguenti funzionalità, relative alla gestione dell'identità digitale nella blockchain di Commerc.io:
	\begin{itemize}
		\item utilizzo e gestione di DID (Decentralized Identifier);
		\item funzioni legate ad un DID: creazione di un DDO (DID Document), richiesta di deposito e richiesta di power up;
		\item gestione degli inviti ad una connessione;
		\item gestione delle credenziali verificabili;
	\end{itemize}
	\item verifica delle funzionalità implementate.
\end{itemize}

\subsection*{Quinta settimana} 
\begin{itemize}
	\item implementazione delle seguenti funzionalità, relative alla gestione e allo scambio di documenti nella blockchain di Commerc.io:
	\begin{itemize}
		\item condivisione di un documento;
		\item invio di una ricevuta;
		\item gestione delle liste di documenti e ricevute;
	\end{itemize}
	\item verifica delle funzionalità implementate.
\end{itemize}

\subsection*{Sesta settimana} 
\begin{itemize}
	\item implementazione dell'interfaccia grafica per la fruizione, da parte degli utenti, delle funzionalità implementate;
	\item validazione del \textit{proof of concept} e verifica dell'integrazione con l'interfaccia.
\end{itemize}

\subsection*{Settima settimana} 
\begin{itemize}
	\item stesura della documentazione relativa a quanto sviluppato;
	\item inizio stesura della relazione sull'attività di stage.
\end{itemize}

\subsection*{Ottava settimana} 
\begin{itemize}
	\item fine stesura della relazione sull'attività di stage;
	\item preparazione della presentazione sull'attività di stage da fare in azienda.
\end{itemize}


\section{}
             					% Stage
% !TEX encoding = UTF-8
% !TEX TS-program = pdflatex
% !TEX root = ../tesi.tex

%**************************************************************
\chapter{Blockchain}
\label{cap:blockchain}

\section{Struttura e funzionamento}

La \gls{blockchaing} è un una tecnologia che si basa su un \textit{database} distribuito che viene acceduto contemporaneamente da più nodi di una rete \textit{peer-to-peer}, ossia una rete di macchine paritetiche che comunicano tra di loro direttamente, senza bisogno di passare attraverso dei nodi aventi delle funzioni di gestione della rete e del traffico presente al suo interno.\\
Attraverso questa rete, i nodi possono effettuare delle transazioni senza il bisogno di un entità terza che faccia da garante intermediario, sfruttando le proprietà della crittografia. Infatti, tramite il suo impiego, le tecnologie \gls{blockchaing} permettono di garantire l'affidabilità sulla veridicità e autenticità delle operazioni e delle comunicazioni effettuate tra individui che non hanno fiducia l'uno dell'altro.\\
In questo contesto, per transazione tra due utenti di una \gls{blockchaing} si intende qualsiasi trasmissione di informazioni, le quali dovranno essere memorizzate all'interno del \textit{database} distribuito, in modo da essere accessibili da parte di tutti i nodi della rete. A questo scopo, le informazioni da trasferire vengono incluse in delle strutture dati chiamate blocchi; questi contengono al loro interno più transazioni e, affinché possano essere aggiunti alla \gls{blockchaing}, ne deve essere verificata la validità e l'affidabilità da parte di tutti i nodi della rete. Per compiere quest'operazione vengono utilizzati opportuni algoritmi di consenso distribuito il cui compito è il raggiungimento di un accordo condiviso su quali blocchi debbano essere aggiunti e quali invece scartati. Una volta che un blocco viene verificato ogni partecipante alla \gls{blockchaing} aggiorna la sua copia locale del \textit{database} con le nuove informazioni.\\
Una \gls{blockchaing} può essere vista come l'insieme di tre parti distinte:

\begin{itemize}
	\item \textbf{\textit{software}:} il \textit{database} distribuito, accessibile solamente in lettura e in inserimento, in cui le informazioni contenute sono verificate come affidabili tramite un protocollo per il consenso distribuito, il quale permette a tutti i nodi di essere d'accordo sull'inserimento o meno di nuove informazioni e, una volta raggiunta l'intesa su un certo dato, gli conferisce lo \textit{status} di valido e non più modificabile; fa parte del \textit{software} anche il livello di rete che realizza lo scambio di informazioni tra i diversi nodi appartenenti alla \gls{blockchaing}
	\item \textbf{crittografia:} impiegata per firmare digitalmente e verificare l'autenticità e l'integrità di tutte le transazioni e i processi di scambio dati tra i nodi della rete;
	\item \textbf{criptovaluta:} una risorsa digitale nativa della \gls{blockchaing}, alla quale viene associato un controvalore in moneta a corso legale, che fornisce gli incentivi economici per l'utilizzo e il mantenimento della rete da parte dei nodi, in modo tale che risulti più conveniente agire in modo onesto piuttosto che in modo disonesto.
\end{itemize}
Attraverso la \gls{blockchaing} è possibile aumentare la trasparenza e la tracciabilità delle filiere produttive e facilitare l'accesso alle transazioni tra privati, aumentandone l'efficienza; il tutto senza bisogno di intermediari che richiedano il pagamento di commissioni e il mantenimento di informazioni personali.

\section{Consenso}

Il protocollo di consenso è un meccanismo resistente agli errori che viene utilizzato all'interno della \gls{blockchaing} per raggiungere l'accordo su un singolo stato della rete tra tutti i nodi distribuiti che la compongono. Questo protocollo è fondamentale per il funzionamento della rete in quanto, dato che ogni nodo può inviare delle nuove transazioni da aggiungere alla \gls{blockchaing}, è necessario che tutte siano sempre controllate e verificate, in modo da raggiungere il consenso su quali debbano essere aggiunte e quali invece rifiutate.\\
Esistono diversi meccanismi di consenso che differiscono per il modo in cui riescono a garantire la fiducia tra i diversi nodi della rete, al fine di raggiungere l'accordo sulle operazioni da eseguire. I principali sono:

\begin{itemize}
	\item \textbf{\gls{pow}:} è stato il primo protocollo di consenso utilizzato ad è conosciuto con il nome di \textit{mining}, per cui i nodi vengono chiamati minatori; questi devono risolvere problemi matematici molto complessi dal punto di vista computazionale, senza però alcun fine pratico, per poter creare un blocco, validare le transazioni in esso contenute e quindi ricevere una ricompensa, in termini di un certo ammontare di criptovaluta, per il lavoro svolto;
	\item \textbf{\gls{pos}:} è un protocollo ideato per essere utilizzato nei contesti in cui i nodi di una rete hanno un proprio interesse nel contribuire al funzionamento della \gls{blockchaing}; in questo scenario viene creato un sottoinsieme di nodi che ha la possibilità di validare i nuovi blocchi e il singolo nodo che avrà il compito di inserire il blocco attuale viene scelto in casualmente tra tutti; per garantire l'affidabilità dei nodi validatori questi sono obbligati a mantenere in stato di fermo un certo ammontare di criptovaluta, in modo da dare prova della loro volontà di agire negli interessi della rete e non nei propri.
\end{itemize}

\section{Immutabilità}

All'interno di una \gls{blockchaing}, ogni volta che una nuova transazione viene inserita, col consenso di tutti i nodi partecipanti alla rete, questa diventa immutabile e non può più essere modificata o eliminata. Infatti, ogni blocco di transazioni all'interno della \gls{blockchaing} sfrutta una funzione crittografica chiamata \textit{hash}, la quale genera una stringa alfanumerica che rappresenta l'intero contenuto del blocco; questa, per le proprietà con cui viene costruita, varia di molto anche per minime modifiche dell'\textit{input} che l'ha generata.\\
Ogni blocco, quindi, contiene al suo interno il proprio valore \textit{hash} e quello del blocco precedente: il primo funge da firma digitale e garantisce l'integrità dei dati su cui è stato calcolato, il secondo invece è necessario all'identificazione del blocco precedente, per poter ripercorrere la \gls{blockchaing} fino all'origine. Sfruttando queste due informazioni si può garantire l'immutabilità in quanto, se venisse fatta una modifica ad un blocco, il suo \textit{hash} verrebbe a sua volta modificato e quindi differirebbe da quello memorizzato all'interno del suo successore, invalidando la catena.

\section{Finalità}

Con finalità si intende il fatto che un blocco validato non può più essere revocato dalla \gls{blockchaing}; in base al protocollo di consenso utilizzato, può essere raggiunta con modalità e in tempi differenti. I principali tipi di finalità sono:

\begin{itemize}
	\item \textbf{finalità probabilistica:} questa tipologia di finalità è caratteristica dei protocolli di tipo \textit{chain-based}, i quali solitamente utilizzando il \gls{pow} come meccanismo di consenso; in queste \gls{blockchaing} la probabilità che un blocco non sia revocato aumenta all'aumentare del numero di blocchi ad esso successivi, quindi più profondo è un blocco più è alta la probabilità che faccia parte della diramazione più lunga della catena, cioè quella valida;
	\item \textbf{finalità istantanea:} questa tipologia di finalità è tipica dei protocolli di tipo \textit{PBFT-based}, i quali solitamente utilizzano il \gls{pos} come meccanismo di consenso; in queste \gls{blockchaing} un blocco è considerato finalizzato non appena viene aggiunto alla catena, in quanto, prima di poter procedere con l'aggiunta, il nuovo blocco deve essere prima proposto da un nodo e poi approvato tramite votazione dalla maggioranza o più dei nodi validatori presenti.
\end{itemize}
             					% Blockchain
% !TEX encoding = UTF-8
% !TEX TS-program = pdflatex
% !TEX root = ../tesi.tex

%**************************************************************
\chapter{Tecnologie e strumenti}
\label{cap:tecnologie-strumenti}

%**************************************************************
\section{Tecnologie}
\subsection*{Commercio.network}

\textit{Commercio.network}\footcite{site:commercio-network} è

\subsection*{Macchina virtuale}

Una macchina virtuale è un \textit{software} che realizza un'emulazione di una macchina fisica attraverso un processo di virtualizzazione, attraverso il quale le risorse \textit{hardware} del dispositivo vengono astratte e rese disponibili al \textit{software} sotto forma virtuale. Di conseguenza non viene fatto accesso e non vengono utilizzate direttamente le risorse fisiche del calcolatore, bensì una loro emulazione logica. Una volta creata una macchina virtuale, su di essa è possibile installare un sistema operativo e delle applicazioni, tramite la stessa procedura che si segue per fare queste operazioni su una macchina fisica.\\
Questa tecnica è usata principalmente per ottimizzare l'utilizzo delle risorse fisiche di un calcolatore, permettendo l'esecuzione contemporanea di più sistemi virtualizzati che svolgono diverse operazioni. Inoltre la virtualizzazione permette di creare sistemi che, essendo slegati dall'\textit{hardware}, non possono in alcun modo causare dei danni al sistema operativo ospitante, anche in caso di corruzione della propria memoria e/o dei propri dischi.\\
Nel progetto di stage è stata utilizzata una macchina virtuale apposita per l'installazione e l'esecuzione di un nodo della rete \textit{Commercio.network}. Così facendo è stato possibile usare la stessa macchina fisica sia per sviluppare e testare il \gls{poc}, all'interno del sistema operativo installato sull'\textit{hardware} reale, sia per interagire con la \gls{blockchaing} di Commerc.io S.r.l., eseguendo tutti i \textit{software} necessari a ciò dentro al sistema virtuale. Solo in questo modo si è potuto mantenere in esecuzione il nodo della rete e, contemporaneamente, avviare ed utilizzare l'applicazione implementata, questa infatti aveva bisogno di accedere alla \gls{blockchaing} per svolgere correttamente tutte le operazioni rese disponibili all'utente.

\subsection*{Git}

Git\footcite{site:git} è

\subsection*{Dart}

Dart\footcite{site:dart} è 

\subsection*{Flutter}

Flutter\footcite{site:flutter} è

%**************************************************************
\section{Strumenti}
\subsection*{VirtualBox}

VirtualBox\footcite{site:virtual-box} è un \textit{software} \gls{open-sourceg} sviluppato da Oracle che fornisce le funzionalità di creazione, controllo ed esecuzione di macchine virtuali ospiti all'interno del sistema operativo del proprio dispositivo fisico. Permette di caricare più sistemi ospiti su un singolo sistema ospitante, ognuno di questi può essere gestito indipendentemente dagli altri, all'interno della propria macchina virtuale. Inoltre, sistema ospitato e sistema ospitante possono comunicare tra di loro sia tramite memoria condivisa che tramite schede di rete virtuali. Infine, supporta sia la virtualizzazione basata su \textit{software} che quella basata su \textit{hardware}, cercando sempre di prediligere quest'ultima alla prima.\\
Nel progetto di stage è stato utilizzato per gestire la macchina virtuale necessaria per l'installazione e l'esecuzione di un nodo della \textit{Commercio.network}.

\subsection*{Ubuntu}

Ubuntu\footcite{site:ubuntu} è un sistema operativo \gls{open-sourceg} basato su Debian, una distribuzione di Linux. Si tratta di un sistema semplice da installare, tramite una procedura guidata, e da utilizzare, attraverso un'intuitiva interfaccia grafica; tuttavia offre molti strumenti utili agli sviluppatori \textit{software}, come gestori di pacchetti, compilatori e ambienti di sviluppo. Inoltre, fa uso di molte tecniche di sicurezza informatica per prevenire il più possibile attacchi sia dall'esterno che dall'interno del sistema. A tal fine prevede una sofisticata gestione dei permessi degli utenti, un controllo totale della connessione e del traffico in rete, e dei compilatori che fanno utilizzo di meccanismi di protezione della memoria. \\
È stato scelto per essere installato all'interno della macchina virtuale utilizzata nel progetto di stage in quanto, tra i sistemi operativi supportati dal \textit{software} che implementa un nodo della rete \textit{Commercio.network}, era quello più adatto ad essere contenuto all'interno dei un ambiente virtuale, sia per la leggerezza in termini di sfruttamento delle risorse, che per la grande capacità di personalizzazione e adattamento.

\subsection*{Cnd}

Cnd è un demone \textit{software}, ossia un programma che, una volta lanciato, resta in esecuzione in \textit{background} senza fornire nessuna possibilità di interazione attiva all'utente. Solitamente questi \textit{software} vengono utilizzati come servizi del sistema operativo e vengono avviati automaticamente durante l'accensione del sistema stesso.\\
Cnd è il \textit{software} sviluppato da Commerc.io S.r.l. per l'implementazione di un nodo della rete \textit{Commercio.network}, è possibile installarlo facilmente seguendo la procedura presente nella documentazione ufficiale\footcite{manual:docs-commercio-network} e, una volta avviato, gestire automaticamente tutte le operazioni necessarie per il corretto funzionamento del nodo e per la corretta interazione con la \gls{blockchaing} in modo da eseguire le operazioni richieste. Oltre a questo, implementa anche un \textit{web server} \gls{rest}\glsfirstoccur{} interrogabile sia dal sistema operativo dove viene eseguito che da dispositivi esterni collegati in rete con esso; è quindi possibile servirsi di questo per operare sulla rete \textit{Commercio.network}.\\
Nel progetto di stage è stato installato all'interno della macchina virtuale Linux per poter poi sfruttarlo per utilizzare la \gls{blockchaing} per la realizzazione delle funzionalità previste per il \gls{poc}. 

\subsection*{Cncli}

Cncli è il software sviluppato da Commerc.io S.r.l. per l'implementazione di un'interfaccia a riga di comando per eseguire operazioni tramite un nodo della rete \textit{Commercio.network} attraverso il demone \textit{software} Cnd. Viene automaticamente installato seguendo la procedura per l'installazione di Cnd e, utilizzandolo dal terminale del sistema operativo, è possibile inserire degli appositi comandi che verranno interpretati e porteranno all'esecuzione di specifiche operazioni sulla \gls{blockchaing} tramite il nodo della rete.\\
Durante il progetto è stato usato, all'interno della macchina virtuale, per testare il funzionamento della \textit{Commercio.network} in base gli scopi dello stage, prima di sviluppare il \gls{poc}. 

\subsection*{Nginx}

Nginx\footcite{site:nginx} è un \textit{software} \gls{open-sourceg} sviluppato da Nginx Inc. che realizza un \textit{web server}. Oltre a questo, può essere anche sfruttato per fare \textit{reverse proxying}, ossia il reindirizzamento delle richieste di un \textit{client} verso uno o più \textit{server} e delle risposte del singolo server verso il \textit{client} opportuno; bilanciamento del carico di richieste dei \textit{client} verso più \textit{server}; \textit{web caching}, ossia la memorizzazione temporanea di pagine \textit{web} recentemente accedute, in modo da reperirle più velocemente per rispondere ad altre richieste successive.\\
Nel progetto di stage è stato utilizzato all'interno della macchina virtuale Linux per esporre all'esterno della stessa le \gls{api} \gls{rest} fornire dal nodo della rete \textit{Commercio.network}; il \textit{server web} eseguito al suo interno, infatti, resta in ascolto per eventuali richieste solamente sull'indirizzo locale del sistema. Tramite la funzionalità di \textit{reverse proxy} di Nginx, quindi, è stato possibile realizzare un reindirizzamento delle richieste fatte all'indirizzo di rete della macchina virtuale, da parte di un dispositivo esterno, verso il \textit{web server} eseguito dal nodo.

\subsection*{GitHub}

GitHub\footcite{site:github} è servizio che permette agli utenti di far ospitare, sulla sua piattaforma \textit{online}, i propri progetti \textit{software}. Fornisce strumenti di sviluppo come il controllo del versionamento distribuito e la gestione del codice sorgente, utilizzando Git, oltre che una serie di altre funzionalità proprie; tra queste si possono trovare sistemi per il controllo d'accesso, per il tracciamento degli errori, per la richiesta di nuove funzionalità e per la gestione dei compiti da svolgere all'interno di un progetto.\\
GitHub offre diversi piani tariffari in base alle funzionalità che si vuole utilizzare, tuttavia, per lo scopo dello stage, è stato sufficiente servirsi del servizio gratuito, in quanto era necessario predisporre un solo progetto privato, senza alcun collaboratore. Infatti la piattaforma è stata utilizzata solamente per effettuare il tracciamento dello sviluppo del codice sorgente del \gls{poc} e il progredire delle sue versioni durante il corso del progetto.

\subsection*{PowerShell}

PowerShell\footcite{site:powershell} è un \gls{frameworkg}\glsfirstoccur{} per l'automatizzazione e la gestione delle configurazioni e dei compiti presenti all'interno del sistema operativo. Consiste di una interfaccia a linea di comando, attraverso la quale l'utente può eseguire diverse operazioni, allo stesso modo che con il terminale di sistema.\\
Nel progetto di stage è stato utilizzato per installare e gestire tutti i \textit{software} necessari per lo sviluppo del \gls{poc}.

\subsection*{Android Studio}

Android Studio\footcite{site:android-studio} è un ambiente di sviluppo software integrato, specificamente progettato per la realizzazione di applicazioni per dispositivi mobili con sistema operativo Android. È stato costruito sulla base di IntelliJ IDEA, un ambiente di sviluppo per i linguaggi Java e Kotlin creato da JetBrains, un'azienda specializzata nella produzione di ambienti di sviluppo integrati per tutti i linguaggi di programmazione più diffusi.\\
Basandosi su IntelliJ, Android Studio ne eredita la maggior parte delle funzionalità, come ad esempio il processo di costruzione dell'applicativo bastato su Gradle, gli strumenti per l'autogenerazione e l'autocorrezione del codice sorgente, e quelli per il rilevamento di problemi di prestazioni, usabilità e compatibilità. Oltre a questo, implementa dei \textit{software} ausiliari per facilitare lo sviluppo delle applicazioni Android, tra questi i più rilevanti sono il generatore di componenti grafici, il visualizzatore dell'anteprima del \textit{layout} dell'interfaccia grafica realizzata e l'emulatore Android che permette di creare un dispositivo virtuale per l'esecuzione e il test delle applicazioni senza necessita di utilizzarne uno fisico.\\
Android Studio è pensato per effettuare delle implementazioni che utilizzano solamente i linguaggi Java o Kotlin, tuttavia mette a disposizione dello sviluppatore una vasta gamma di espansioni che ne espandono le funzionalità. Per il progetto di stage, infatti, è stata prevista la realizzazione di un \gls{poc} in linguaggio Dart tramite il \gls{frameworkg} Flutter; a tal fine sono state installati gli opportuni pacchetti per consentire l'estensione delle funzioni base dell'ambiente di sviluppo anche al codice scritto in Dart/Flutter.
             					% Tecnologie e strumenti
% !TEX encoding = UTF-8
% !TEX TS-program = pdflatex
% !TEX root = ../tesi.tex

%**************************************************************
\chapter{Applicazione}
\label{cap:applicazione}

%**************************************************************
\section{Analisi}

\subsection{Descrizione}

Durante la prima parte dello stage, svolgendo le attività di studio e formazione personale sul contesto aziendale e sulla \gls{blockchaing} sviluppata da Commerc.io S.r.l., sono stati individuati i temi di principale interesse per \myCompany{} \companyTitle:

\begin{itemize}
	\item la gestione delle identità digitali e delle informazioni che le caratterizzano;
	\item lo scambio di documenti con valore legale associati ad un'identità tramite la relativa firma elettronica.
\end{itemize}

Di conseguenza, l'obiettivo principale della seconda fase del progetto è stato realizzare un prototipo che dimostrasse la fattibilità dell'implementazione di un'applicazione per dispositivi mobili inerente alle due precedenti tematiche. Questa doveva permettere agli utenti di registrare una propria \gls{ssi} e, una volta completata l'apposita procedura, di utilizzarla per certificare la paternità degli eventuali documenti caricati e poi condivisi con altri utenti, tramite il programma stesso.\\
L'applicazione dovrà essere disponibile sia per dispositivi Android che per dispositivi iOS, mantenendo un'interfaccia grafica il più coerente e simile possibile, indipendentemente dal sistema operativo su cui veniva eseguita, e fornendo le stesse funzionalità. In particolare, ogni utente dovrà poter:

\begin{itemize}
	\item registrare una nuova identità, la quale sarà poi accessibile ed utilizzabile solamente all'interno del sistema attualmente utilizzato;
	\item recuperare un'identità precedentemente creata in un dispositivo differente da quello corrente;
	\item accedere alla propria identità in modo sicuro;
	\item utilizzare l'identità all'interno dell'applicazione per effettuare le diverse operazioni messe a disposizione.
\end{itemize}

Le prime tre funzionalità richiederanno l'esecuzione di alcune operazioni che, poste in una sequenza ordinata, andranno a formare una procedura che dovrà essere seguita passo passo per poter raggiungerne lo scopo.\\\\
Per registrare una nuova identità sarà necessario inserire il proprio nome e cognome e scegliere il metodo di autenticazione desiderato tra quelli disponibili, ossia \textit{password} alfanumerica o dati biometrici; nel caso venga fatta la prima scelta verrà richiesto l'inserimento, con successiva conferma, della \textit{password} desiderata, altrimenti sarà necessario autenticarsi con l'impronta digitale o con l'identificativo facciale prima di poter proseguire. Successivamente verrà chiesto di inserire il proprio indirizzo di posta elettronica e, una volta confermato, un codice di verifica numerico inviato all'utente tramite \textit{email} per certificare l'effettiva proprietà della casella indicata. La stessa operazione dovrà essere poi ripetuta per l'inserimento del numero di telefono cellulare, in questo caso la verifica verrà effettuata tramite \gls{sms}\glsfirstoccur. Infine verrà generato e visualizzato a schermo il codice mnemonico usato per generare il \gls{walletg} associato alla \gls{ssi} appena creata, e sarà possibile salvarlo in formato \gls{pdf}\glsfirstoccur{} all'interno della memoria del dispositivo, oppure stamparlo direttamente dall'applicazione. È di vitale importanza far capire all'utente che non deve smarrire questo codice, in quanto è l'unica informazione utilizzabile per ripristinare l'accesso alla propria identità.\\\\
Per poter recuperare una vecchia identità ed accedervi con un nuovo dispositivo, diverso da quello utilizzato per crearla la prima volta, verrà prima richiesto il metodo di autenticazione desiderato, allo stesso modo della procedura di registrazione e, dopo aver fornito i dati necessari, dovrà essere inserito il codice mnemonico generato durante la creazione della \gls{ssi}. Solo nel caso questo sia corretto l'utente potrà accedere alla propria identità, altrimenti la procedura terminerà con un'errore e, per poter utilizzare l'applicazione, sarà quindi indispensabile ripetere nuovamente il procedimento di registrazione dall'inizio.\\\\
Per effettuare l'accesso, invece, non sarà necessario inserire il codice mnemonico, ma sarà sufficiente fornire i dati di autenticazione richiesti in base al metodo scelto in precedenza. Se era stato selezionato il riconoscimento biometrico, allora l'utente dovrà utilizzare la propria impronta digitale o il riconoscimento facciale per entrare, altrimenti verrà richiesto l'inserimento della \textit{password} alfanumerica scelta dall'utente durante la registrazione.\\\\
Una volta all'interno dell'applicazione, le funzionalità che saranno a disposizione dell'utente sono:

\begin{itemize}
	\item visualizzazione della foto e delle informazioni associate alla propria identità, corredate dall'indirizzo pubblico e dal bilancio del \gls{walletg} ad essa associato;
	\item visualizzazione della foto, del nome e delle informazioni associate ad ogni documento inserito all'interno dell'applicazione;
	\item visualizzazione del codice QR associato all'indirizzo pubblico del proprio \gls{walletg};
	\item inserimento di una nuova foto profilo, la quale può essere scattata direttamente all'interno dell'applicazione;
	\item inserimento di un nuovo documento.
\end{itemize}

\subsection{Requisiti}

L'obiettivo principale che si desidera raggiungere è l'implementazione di un \gls{poc} che dimostri la fattibilità dell'applicazione d'interesse per l'azienda, di conseguenza il progetto non prevede uno sviluppo \textit{software} completo, ma limitato a rendere disponibili le operazioni più significative. Quindi, dopo un'attenta analisi del problema proposto e delle funzionalità richieste per l'applicazione, tenendo anche conto del tempo rimasto a disposizione per il progetto di stage, sono stati individuati diversi tipi di requisiti, in base all'importanza e alla quantità di lavoro richiesta per riuscire a soddisfarli.\\
Per identificare ogni requisito in modo univoco si è fatto utilizzo di un opportuno codice, il quale è strutturato nel modo seguente:\\

\begin{center}
	\textbf{R[T][P]-[N]}
\end{center}

Dove ogni lettera ha un preciso significato:

\begin{itemize}
	\item \textbf{R:} requisito;
	
	\item \textbf{T:} tipologia, differisce in base allo specifico requisito e può assumere i seguenti valori:
		\begin{itemize}
			\item \textbf{F:} funzionale, ossia un requisito che indica una caratteristica operativa che il \textit{software} deve avere;
			\item \textbf{Q:} qualitativo, ossia un requisito che indica una caratteristica, non necessariamente legata al \textit{software}, ma che ne aumenta la qualità in termini di efficacia ed efficienza;
			\item \textbf{V:} vincolo, ossia un requisito che indica una caratteristica, non necessariamente legata al \textit{software}, ma che deve essere garantita su richiesta dell'azienda;
		\end{itemize}
	
	\item \textbf{P:} priorità, differisce in base allo specifico requisito e può assumere i seguenti valori:
		\begin{itemize}
			\item \textbf{O:} obbligatorio, ossia un requisito vincolante in quanto primario e fondamentale;
			\item \textbf{D:} desiderabile, ossia un requisito non vincolante o strettamente necessario, ma dal riconoscibile valore aggiunto;
			\item \textbf{F:} funzionale, ossia un requisito rappresentante valore aggiunto non strettamente competitivo; 
		\end{itemize}
	
	\item \textbf{N:} numero, intero progressivo maggiore di zero con funzione di identificativo del requisito.
\end{itemize}

Di seguito vengono riportati tutti i requisiti riguardanti lo sviluppo del \gls{poc}, suddivisi in tre tabelle, una per ogni tipologia.

\rowcolors{1}{grayer}{white}
\begin{longtable}{|c|p{10.5cm}|}
	\hline
	\rowcolor{gray}
	\textbf{Requisito} & \textbf{Descrizione} \\
	\hline
	RFO-1     & L'applicazione deve permettere l'inserimento del nome dell'utente. \\
	\hline
	RFO-2     & L'applicazione deve permettere l'inserimento del cognome dell'utente. \\
	\hline
	RFO-3     & L'applicazione deve permettere l'inserimento dell'indirizzo di posta elettronica dell'utente. \\
	\hline
	RFO-4     & L'applicazione deve permettere l'inserimento del numero di telefono cellulare dell'utente. \\
	\hline
	RFO-5     & L'applicazione deve permettere l'inserimento del codice di verifica dell'indirizzo di posta elettronica dell'utente. \\
	\hline
	RFO-6     & L'applicazione deve permettere l'inserimento del codice di verifica del numero di telefono cellulare dell'utente. \\
	\hline
	RFO-7     & L'applicazione deve permettere l'invio del codice di verifica dell'indirizzo di posta elettronica dell'utente. \\
	\hline
	RFO-8     & L'applicazione deve permettere l'invio del codice di verifica del numero di telefono cellulare dell'utente. \\
	\hline
	RFO-9     & L'applicazione deve permettere la visualizzazione del codice mnemonico. \\
	\hline
	RFO-10    & L'applicazione deve permettere l'inserimento del codice mnemonico. \\
	\hline
	RFO-11    & L'applicazione deve permettere l'inserimento di una \textit{password} per l'accesso da parte dell'utente. \\
	\hline
	RFO-12    & L'applicazione deve permettere l'inserimento della \textit{password} di accesso dell'utente. \\
	\hline
	RFO-13    & L'applicazione deve permettere la scelta dell'accesso tramite \textit{password} da parte dell'utente. \\
	\hline
	RFO-14    & L'applicazione deve permettere l'inserimento della password per l'accesso da parte dell'utente. \\
	\hline
	RFO-15    & L'applicazione deve permettere la visualizzazione del nome dell'utente. \\
	\hline
	RFO-16    & L'applicazione deve permettere la visualizzazione del cognome dell'utente. \\
	\hline
	RFO-17    & L'applicazione deve permettere la visualizzazione dell'indirizzo di posta elettronica dell'utente. \\
	\hline
	RFO-18    & L'applicazione deve permettere la visualizzazione del numero di telefono cellulare dell'utente. \\
	\hline
	RFO-19    & L'applicazione deve permettere la visualizzazione dell'indirizzo pubblico del \gls{walletg} dell'utente. \\
	\hline
	RFO-20    & L'applicazione deve permettere la visualizzazione del saldo del \gls{walletg} dell'utente. \\
	\hline
	RFO-21    & L'applicazione deve permettere la visualizzazione della foto profilo dell'utente. \\
	\hline
	RFO-22    & L'applicazione deve permettere la visualizzazione della lista dei documenti inseriti dall'utente. \\
	\hline
	RFO-23    & L'applicazione deve permettere la registrazione di una nuova identità da parte dell'utente. \\
	\hline
	RFO-24    & L'applicazione deve permettere il recupero di un'identità creata in precedenza da parte dell'utente. \\
	\hline
	RFO-25    & L'applicazione deve permettere l'accesso tramite autenticazione da parte dell'utente. \\
	\hline
	RFD-1     & L'applicazione deve permettere il reinvio del codice di verifica dell'indirizzo di posta elettronica dell'utente. \\
	\hline
	RFD-2     & L'applicazione deve permettere il reinvio del codice di verifica del numero di telefono cellulare dell'utente. \\
	\hline
	RFD-3     & L'applicazione deve permettere il salvataggio del codice mnemonico in formato \gls{pdf}. \\
	\hline
	RFD-4     & L'applicazione deve permettere la scelta dell'accesso tramite dati biometrici da parte dell'utente. \\
	\hline
	RFD-5     & L'applicazione deve permettere l'inserimento dei dati biometrici per l'accesso da parte dell'utente. \\
	\hline
	RFD-6     & L'applicazione deve permettere l'inserimento della foto profilo da parte dell'utente. \\
	\hline
	RFD-7     & L'applicazione deve permettere l'inserimento del titolo di un documento da parte dell'utente. \\
	\hline
	RFD-8     & L'applicazione deve permettere l'inserimento delle informazioni di un documento da parte dell'utente. \\
	\hline
	RFD-9     & L'applicazione deve permettere l'inserimento di un documento da parte dell'utente. \\
	\hline
	RFF-1     & L'applicazione deve permettere la stampa del codice mnemonico. \\
	\hline
	RFF-2     & L'applicazione deve permettere la visualizzazione del codice QR associato all'indirizzo pubblico del \gls{walletg}. \\
	\hline
	RFF-3     & L'applicazione deve permettere la modifica della foto profilo da parte dell'utente. \\
	\hline
	RFF-4     & L'applicazione deve permettere l'inserimento della foto di un documento da parte dell'utente. \\
	\hline
	
	\caption{Tabella dei requisti funzionali}
	\label{tab:requisiti-funzionali}
\end{longtable}

\rowcolors{1}{white}{grayer}
\begin{longtable}{|c|p{10.5cm}|}
	\hline
	\rowcolor{gray}
	\textbf{Requisito} & \textbf{Descrizione} \\
	\hline
	RQO-1    & Deve essere fornita una documentazione delle procedure di installazione e configurazione di tutte le tecnologie e gli strumenti utilizzati per sviluppare e testare l'applicazione. \\
	\hline
	RQO-2    & Deve essere fornita una documentazione di tutte le classi che compongono l'applicazione, con i relativi metodi. \\
	\hline
	RQO-3    & Il processo di sviluppo deve fare uso di un sistema di versionamento. \\
	\hline
	
	\caption{Tabella dei requisiti qualitativi}
	\label{tab:requisiti-qualitativi}
\end{longtable}

\rowcolors{1}{white}{grayer}
\begin{longtable}{|c|p{10.5cm}|}
	\hline
	\rowcolor{gray}
	\textbf{Requisito} & \textbf{Descrizione} \\
	\hline
	RVO-1    & L'applicazione deve essere scritta utilizzando il linguaggio Dart. \\
	\hline
	RVO-2    & L'applicazione deve essere scritta utilizzando il \gls{frameworkg} Flutter. \\
	\hline
	RVO-3    & L'applicazione deve interagire con la \gls{blockchaing} della \textit{Commercio.network}. \\
	\hline
	RVO-4    & L'interazione a basso livello con la \gls{blockchaing} deve avvenire tramite la libreria Sacco, sviluppata da Commerc.io S.r.l. per il linguaggio di programmazione Dart. \\
	\hline
	RVO-5    & L'interazione ad alto livello con la \gls{blockchaing} deve avvenire tramite la libreria CommercioSDK, sviluppata da Commerc.io S.r.l. per il linguaggio di programmazione Dart. \\
	\hline
	RVO-6    & L'applicazione deve essere eseguibile su sistema operativo Android. \\
	\hline
	RVD-1    & L'applicazione deve essere eseguibile su sistema operativo iOS. \\
	\hline
	
	\caption{Tabella dei requisiti di vincolo}
	\label{tab:requisiti-vincolo}
\end{longtable}
	

%**************************************************************
\section{Progettazione}

Per lo studio dei casi di utilizzo del prodotto sono stati creati dei diagrammi.\\
I diagrammi dei casi d'uso (in inglese \emph{Use Case Diagram}) sono diagrammi di tipo \gls{uml} dedicati alla descrizione delle funzioni o servizi offerti da un sistema, così come sono percepiti e utilizzati dagli attori che interagiscono col sistema stesso.\\
Essendo il progetto finalizzato alla creazione di un tool per l'automazione di un processo, le interazioni da parte dell'utilizzatore devono essere ovviamente ridotte allo stretto necessario. Per questo motivo i diagrammi d'uso risultano semplici e in numero ridotto.

\subsection*{Namespace 1}
Descrizione namespace 1.

\begin{namespacedesc}
    \classdesc{Classe 1}{Descrizione classe 1}
    \classdesc{Classe 2}{Descrizione classe 2}
\end{namespacedesc}

%**************************************************************
\section{Codifica}

%**************************************************************
\section{Possibili evoluzioni}
             					% Applicazione
% !TEX encoding = UTF-8
% !TEX TS-program = pdflatex
% !TEX root = ../tesi.tex

%**************************************************************
\chapter{Conclusioni}
\label{cap:conclusioni}

%**************************************************************
\section{Consuntivo finale}

%**************************************************************
\section{Raggiungimento degli obiettivi}

%**************************************************************
\section{Conoscenze e competenze acquisite}

%**************************************************************
\section{Valutazione personale}
             					% Conclusioni
%% !TEX encoding = UTF-8
% !TEX TS-program = pdflatex
% !TEX root = ../tesi.tex

%**************************************************************
\chapter{Conclusioni}
\label{cap:}
%**************************************************************

%**************************************************************
\section{Consuntivo finale}

%**************************************************************
\section{Raggiungimento degli obiettivi}

%**************************************************************
\section{Conoscenze acquisite}

%**************************************************************
\section{Valutazione personale}

%\appendix                               
%\input{capitoli/capitolo-A}

%**************************************************************
% Materiale finale
%**************************************************************

\backmatter
\printglossary[type=acronym, title={Acronimi e abbreviazioni}]
\printglossary[type=main, title={Glossario}]
% !TEX encoding = UTF-8
% !TEX TS-program = pdflatex
% !TEX root = ../tesi.tex

%**************************************************************
% Bibliografia
%**************************************************************

\cleardoublepage
\chapter{Bibliografia}

\nocite{*}
% Stampa i riferimenti bibliografici
\printbibliography[heading=subbibliography, title={Riferimenti bibliografici}, type=book]

% Stampa i siti web consultati
\printbibliography[heading=subbibliography, title={Siti web consultati}, type=online]


\end{document}
